\section{Conclusion}
\label{sec:conclusion}

We set out to test whether measuring a simulator's mismatches against the
real system makes its closed-loop predictions transfer, and we committed
the predictions before going to the water so that the test would be a
test. The stack did not transfer, and the two differences that shaped the
physical experiment most were not calibrated quantities at all: the
detector's acceptance gates, and the size of the room.

The honest reading is that our calibration effort was aimed at the
quantities that were easy to measure rather than the ones that governed
the outcome. Detection probability, range bias, dropout burstiness and
actuator response are all real, all were measured against the physical
vehicle, and all were faithfully injected; none of them determined
whether the system worked in the pool.

The post-deployment ablation sharpened that reading in two ways we did
not expect. The largest single effect on simulated safety was not a
calibration level or a field change but a defect in our own frozen
simulator: correcting the relay's field of view raised median clearance
from \SI{1.29}{\metre} to \SI{2.91}{\metre} and removed its collisions.
And giving the planner a manoeuvre that the basin can actually contain
collapses clearance to \SI{0.30}{\metre} with four collisions in ten
runs, which makes the pool an unfair test of this planner rather than a
difficult one: its safety margin is the width of its detour, and the
basin has no room for it.

We think four things generalise. First, the acceptance gates of a
perception module belong in its published interface, because they define
the range interval over which it produces any observation at all, and
that interval can silently exclude the experiment. Second, a sim-to-real
study should derive its scenario geometry from the physical site before
freezing anything, since a manoeuvre the site cannot contain turns a
controller result into an architecture result. Third, the boundaries of
the physical workspace belong in the simulated scenario: ours had none,
which is why every simulated run completed a manoeuvre that no physical
run could, and no calibration of sensing or actuation could have
supplied that. Fourth, and most uncomfortably, a campaign that adapts a
stack in the field has stopped testing the stack it froze; saying so
plainly costs a claim and buys back the reader's ability to judge
everything else.

What the physical campaign does support is narrow and solid: a
camera-driven avoidance controller ran closed-loop on a BlueROV2 and
avoided a real obstacle in eight of eight admissible runs, at a commanded
surge independent of geometry, choosing its side from where it saw the
obstacle rather than from a fixed preference -- the last of these
established by mirroring the geometry mid-session, with two runs, and
reported as the exploratory observation it is. What it cannot support,
and what we therefore do not claim, is any causal validation of the
calibration ladder, any ranking of the two planners in water, and
anything at all about the vehicle's trajectory, which was never
measured.

The pool is drained and the physical dataset is closed. The
post-deployment simulations reported here were run afterwards, with the
outcome known, and are labelled as such throughout; they identify which
of the deployment changes carried the effect, which a single day in the
water could not. We would rather publish that separation, together with
an audit of our own frozen study's defects, than a cleaner-looking result
that quietly conflates a prediction with a diagnosis.
